%% 31장 스플라인
\chapter{Splines}
\label{ch:splines}

\abstract{앞의 두 장은 함수 전체를 하나의 다항식이나 유리함수로 근사했다. 이 장에서는 정반대의 길을 간다. 구간을 잘게 나누어 조각마다 낮은 차수의 다항식을 두고 조각들을 경계에서 매끄럽게 이어 붙인다. 이런 조각 다항식을 스플라인이라 한다. 노드의 분포를 고를 필요가 없고 차수가 높아지지 않으므로 고른 간격의 자료에서도 안정하다. 조각 3차 Hermite 보간에서 출발해 이음매의 2차 도함수까지 이어지게 하는 3차 스플라인이 삼중대각 계로 정해진다는 것과 그 계를 피벗 없이 풀어도 안전한 이유를 증명한다. 3차 스플라인이 곡률 에너지를 최소화한다는 고전적인 정리, 경계 조건에 따라 수렴 차수가 달라지는 현상, 단조로운 자료에서 넘침 없이 보간하는 PCHIP를 차례로 본다. 마지막으로 스플라인 공간의 국소적 기저인 B-스플라인과 그 평가 알고리즘, 보간 행렬의 완전 비음성을 다룬다.}

\section{Piecewise Polynomials}
\label{sec:ppoly}

\subsection{Representation}

\term{중단점}{breakpoint} $x_0<x_1<\dots<x_n$으로 나눈 구간마다 차수 $k$ 이하의 다항식을 두는 함수를 \term{조각 다항식}{piecewise polynomial}이라 한다. 컴퓨터에서는 흔히 조각마다 그 조각의 왼쪽 끝을 중심으로 한 거듭제곱 기저의 계수를 저장한다.
\[
p(x)=\sum_{e=0}^kc_{e,i}(x-x_i)^e,\qquad x_i\le x<x_{i+1}
\]
값을 구하려면 먼저 $x$가 속한 조각 $i$를 15장의 이분 탐색으로 찾고, $s=x-x_i$에서 앞 장들의 Horner 방법으로 평가한다. 조각의 길이가 짧으므로 $s$가 작고, 29장의 단항식 기저의 문제는 생기지 않는다. 조각 다항식의 도함수와 역도함수도 조각마다 계수를 바꾸면 되므로 쉽다. 역도함수는 중단점에서 값이 이어지도록 조각마다 상수를 누적해야 한다.

조각 다항식으로 매끄러운 함수를 보간하면 각 조각의 오차는 29장의 \cref{thm:interp-error}에 따라 조각 길이 $h$의 $k+1$제곱에 비례한다. 차수를 올리는 대신 조각을 잘게 나누어 정확도를 얻으므로 Runge 현상이 일어날 여지가 없다. 대신 수렴은 기하급수적이지 않고 $h^{k+1}$의 대수적 속도에 머문다.

\subsection{Piecewise Cubic Hermite Interpolation}

가장 많이 쓰이는 것은 $k=3$이다. 중단점에서 값 $y_i$와 기울기 $s_i$가 주어지면 조각마다 양 끝의 값과 기울기 네 개를 맞추는 3차 다항식이 하나로 정해진다. 이것이 29장의 Hermite 보간의 가장 단순한 경우이다. $h_i=x_{i+1}-x_i$, $\Delta_i=(y_{i+1}-y_i)/h_i$로 쓰면 그 계수는
\begin{equation}
\label{eq:hermite-coef}
c_0=y_i,\quad c_1=s_i,\quad c_2=\frac{3\Delta_i-2s_i-s_{i+1}}{h_i},\quad c_3=\frac{s_i+s_{i+1}-2\Delta_i}{h_i^2}
\end{equation}
이다. $p(x_{i+1})=y_{i+1}$과 $p'(x_{i+1})=s_{i+1}$을 대입해 확인할 수 있다. 이렇게 만든 함수는 값과 1차 도함수가 이음매에서 이어지는 $C^1$ 함수이다. 기울기가 $f$의 정확한 도함수이면 오차는 다음과 같다.

\begin{proposition}
\label{prop:hermite-error}
$f\in C^4$이고 $s_i=f'(x_i)$이면 조각 $[x_i,x_{i+1}]$에서 $|f(x)-p(x)|\le h_i^4\max|f^{(4)}|/384$이다.
\end{proposition}

\begin{proof}
29장의 \cref{thm:interp-error}는 겹친 노드에도 그대로 성립한다(증명의 Rolle 논증에서 겹친 근을 중복도대로 세면 된다). 노드 다항식이 $\omega(x)=(x-x_i)^2(x-x_{i+1})^2$이고 그 최댓값은 가운데에서 $(h_i/2)^4=h_i^4/16$이므로, 오차는 $\max|f^{(4)}|/4!\cdot h_i^4/16$ 이하이다.
\end{proof}

그러나 실제 자료에는 기울기가 없다. 기울기를 어떻게 정하느냐에 따라 아래의 3차 스플라인, PCHIP, Akima 보간이 갈라진다. stardust를 비롯한 여러 라이브러리는 이 셋을 모두 같은 틀로 구현한다. 기울기를 먼저 구하고 그 기울기로 \eqref{eq:hermite-coef}의 계수를 만든다.

\section{Cubic Splines}
\label{sec:cubic-spline}

\subsection{The Tridiagonal System}

기울기를 자유롭게 고를 수 있다면 그 자유를 써서 2차 도함수까지 이음매에서 이어지게 할 수 있다. 그렇게 만든 $C^2$ 조각 3차 다항식이 \term{3차 스플라인}{cubic spline}이다. \eqref{eq:hermite-coef}에서 조각 $i$의 왼쪽 끝 2차 도함수는 $2c_2=2(3\Delta_i-2s_i-s_{i+1})/h_i$이고, 조각 $i-1$의 오른쪽 끝 2차 도함수는 $2c_2+6c_3h$를 정리한 $2(-3\Delta_{i-1}+s_{i-1}+2s_i)/h_{i-1}$이다. 두 값이 같다는 조건을 정리하면 안쪽 노드 $i=1,\dots,n-1$마다
\begin{equation}
\label{eq:spline-system}
h_is_{i-1}+2(h_{i-1}+h_i)s_i+h_{i-1}s_{i+1}=3(h_i\Delta_{i-1}+h_{i-1}\Delta_i)
\end{equation}
를 얻는다. 미지수는 $n+1$개이고 식은 $n-1$개이므로 양 끝에서 조건이 하나씩 더 필요하다. 흔히 쓰는 것은 다음 넷이다.
\begin{enumerate}
\item \term{자연 스플라인}{natural spline}: 양 끝의 2차 도함수가 $0$이다. 첫 행은 $2h_0s_0+h_0s_1=3h_0\Delta_0$이다.
\item \term{고정 스플라인}{clamped spline}: 양 끝의 기울기 $s_0$, $s_n$을 준다.
\item \term{not-a-knot}{not-a-knot}: 첫 두 조각과 마지막 두 조각이 각각 같은 3차 다항식이다. 곧 $x_1$과 $x_{n-1}$에서 3차 도함수까지 이어진다.
\item \term{주기 스플라인}{periodic spline}: $y_0=y_n$이고 양 끝의 1차, 2차 도함수가 같다. 계가 순환 삼중대각이 된다.
\end{enumerate}

자연 스플라인과 고정 스플라인의 계수 행렬은 각 행에서 대각이 비대각 원소의 합의 두 배 이상이다. 안쪽 행은 $2(h_{i-1}+h_i)$ 대 $h_i+h_{i-1}$이고, 끝 행은 $2h_0$ 대 $h_0$이거나 $1$ 대 $0$이다. 따라서 4장의 Gershgorin 정리로 행렬은 정칙이고 스플라인은 유일하다. 이 행렬을 피벗 없이 소거해도 안전하다.

\begin{proposition}
\label{prop:thomas-positive}
자연 스플라인과 고정 스플라인의 계를 피벗 없이 소거하는 \term{Thomas 알고리즘}{Thomas algorithm}에서, 모든 피벗 $u_i$는 그 행의 위 대각 원소 $c_i$ 이상인 양수이고, 소거의 승수와 $L$, $U$의 원소는 모두 $0$ 이상이다.
\end{proposition}

\begin{proof}
삼중대각 행렬의 행 $i$를 $[a_i,b_i,c_i]$로 쓰면 소거는 $l_i=a_i/u_{i-1}$, $u_i=b_i-l_ic_{i-1}$이다. 모든 원소가 $0$ 이상이다. 첫 행에서 $u_0=b_0\ge c_0$이다($2h_0\ge h_0$ 또는 $1\ge0$). $u_{i-1}\ge c_{i-1}$이라 하면 $l_ic_{i-1}=a_ic_{i-1}/u_{i-1}\le a_i$이므로 안쪽 행에서는 $u_i\ge2(h_{i-1}+h_i)-h_i=2h_{i-1}+h_i\ge h_{i-1}=c_i$이다. 마지막 행도 같은 계산으로 양수이다. 승수 $l_i$는 음이 아닌 수의 몫이므로 $0$ 이상이다.
\end{proof}

$L$과 $U$가 모두 비음이면 $|L||U|=|LU|=|A|$이므로, 9장과 19장에서 본 LU의 후진 오차 상계 $|\Delta A|\le\gamma_k|\hat L||\hat U|$는 $|A|$의 상수배가 된다. 곧 피벗 없는 Thomas 알고리즘은 성분별로 후진 안정하다. 삼중대각 계의 풀이는 $O(n)$의 일이다. 자료가 여러 열이면 분해를 한 번 하고 대입만 열마다 하면 된다.

\subsection{Minimum Curvature}

3차 스플라인이라는 이름은 제도사들이 곡선을 그릴 때 쓰던 휘어지는 자에서 왔다. 가는 나무 자를 핀으로 점들에 고정하면 자는 탄성 에너지를 최소로 하는 모양이 되는데 작은 변형에서 그 에너지는 대략 $\int(y'')^2$에 비례한다. 자연 스플라인은 정확히 이 에너지를 최소화한다\cite{holladay1957}.

\begin{theorem}[Holladay]
\label{thm:holladay}
$g\in C^2[x_0,x_n]$이 모든 노드에서 $g(x_i)=y_i$이면, 같은 자료의 자연 스플라인 $s$에 대해
\[
\int_{x_0}^{x_n}s''(x)^2\,dx\le\int_{x_0}^{x_n}g''(x)^2\,dx
\]
이고, 등호는 $g=s$일 때뿐이다.
\end{theorem}

\begin{proof}
$e=g-s$로 두면 $\int g''^2=\int s''^2+2\int s''e''+\int e''^2$이므로 $\int s''e''=0$을 보이면 된다. 조각마다 부분적분하면
\[
\int_{x_i}^{x_{i+1}}s''e''\,dx=\bigl[s''e'\bigr]_{x_i}^{x_{i+1}}-\int_{x_i}^{x_{i+1}}s'''e'\,dx
\]
이다. $s''$은 연속이므로 첫 항을 조각에 걸쳐 더하면 양 끝의 값만 남는데, 자연 스플라인은 $s''(x_0)=s''(x_n)=0$이므로 $0$이다. 둘째 항에서 $s'''$은 조각마다 상수 $\sigma_i$이므로 $\sigma_i\int_{x_i}^{x_{i+1}}e'=\sigma_i\bigl(e(x_{i+1})-e(x_i)\bigr)=0$이다. 노드에서 $e=0$이기 때문이다. 등호가 성립하면 $\int e''^2=0$이므로 $e$는 일차식이고, 노드에서 $0$이므로 $e=0$이다.
\end{proof}

같은 자료에서 기울기 조건이 다른 보간과 견주면 정리를 수치로 확인할 수 있다. $[0,1]$의 무작위 점 $12$개에 잡음 섞인 값을 두고 곡률 에너지 $\int s''^2$을 재면, 자연 스플라인이 $56683$, not-a-knot이 $58423$, 양 끝 기울기를 $0$으로 고정한 스플라인이 $58945$이고, 아래의 PCHIP는 $211978$이다.

\subsection{Boundary Conditions and Order of Convergence}

자연 스플라인은 에너지가 가장 작지만 근사로서는 대가가 있다. 실제 함수의 2차 도함수가 양 끝에서 $0$이 아니면 그 조건이 거짓이기 때문이다. $[0,1]$에서 $e^x$를 고른 간격으로 보간하면 \cref{tab:spline-bc}와 같다.

\begin{table}[t]
\begin{sdtable}{colspec={cccc}, row{1}={font=\small\bfseries}}
\toprule
$n$ & 자연 & 고정 & not-a-knot \\
\midrule
$10$ & $1.33\times10^{-3}$ & $6.96\times10^{-7}$ & $6.93\times10^{-6}$ \\
$20$ & $3.34\times10^{-4}$ & $4.39\times10^{-8}$ & $4.56\times10^{-7}$ \\
$40$ & $8.34\times10^{-5}$ & $2.75\times10^{-9}$ & $2.92\times10^{-8}$ \\
$80$ & $2.09\times10^{-5}$ & $1.72\times10^{-10}$ & $1.85\times10^{-9}$ \\
$160$ & $5.21\times10^{-6}$ & $1.08\times10^{-11}$ & $1.16\times10^{-10}$ \\
\bottomrule
\end{sdtable}
\caption{$[0,1]$의 $e^x$를 고른 간격 $n$구간의 3차 스플라인으로 보간했을 때의 최대 오차. 고정 스플라인은 정확한 끝 기울기 $1$, $e$를 썼다.}
\label{tab:spline-bc}
\end{table}

고정 스플라인과 not-a-knot은 $n$이 두 배가 될 때마다 오차가 $16$배 줄어드는 $O(h^4)$이고, 자연 스플라인은 $4$배씩 줄어드는 $O(h^2)$이다. 자연 스플라인의 오차는 양 끝 근처에 몰려 있다. 고정 스플라인에 대해서는 Hall과 Meyer가 최적의 상수까지 구했다. $\|f-s\|_\infty\le(5/384)h^4\|f^{(4)}\|_\infty$이다\cite{hall1976}. $n=10$에서 이 상계는 $3.5\times10^{-6}$이고 실제 오차는 $7.0\times10^{-7}$이다. 그래서 정확한 끝 기울기를 모를 때는 not-a-knot을 표준으로 쓴다. 끝에 거짓 정보를 넣지 않으면서도 $O(h^4)$을 지키기 때문이다.

not-a-knot의 끝 행은 대각 우세가 아니다. 첫 행이 $[h_1,\ x_2-x_0]$으로 대각보다 비대각이 크다. 그래도 첫 피벗이 $h_1$, 다음 승수가 정확히 $1$이 되어 둘째 피벗이 $h_0+h_1\ge c_1$로 회복되므로 \cref{prop:thomas-positive}의 논증이 둘째 행부터 다시 성립한다. 마지막 행 $[x_n-x_{n-2},\ h_{n-2}]$도 대각 우세가 아니지만, 바로 앞의 피벗이 $u_{n-1}\ge2h_{n-2}+h_{n-1}$이므로 마지막 피벗은 $h_{n-2}\bigl(1-(h_{n-2}+h_{n-1})/(2h_{n-2}+h_{n-1})\bigr)>0$ 이상이다.

\section{Shape-Preserving Interpolation}
\label{sec:shape-preserving}

\subsection{Overshoot}

3차 스플라인은 매끄럽지만 자료가 갑자기 변하는 곳에서 자료의 범위를 넘어 출렁인다. $x=0,\dots,7$에서 값이 $0,0,0,0.1,0.9,1,1,1$인 계단 모양의 자료를 자연 스플라인으로 보간하면 곡선이 $-0.0284$까지 내려갔다가 $1.0284$까지 올라가고, not-a-knot도 $\pm0.0291$만큼 넘친다. 농도나 확률처럼 음수가 될 수 없는 양이나 단조로워야 하는 누적 분포를 보간할 때 이런 넘침은 받아들일 수 없다. 2차 도함수의 연속성을 포기하고 기울기를 국소적으로 정하면 이 문제를 피할 수 있다.

\subsection{Conditions for Monotonicity}

단조 증가하는 자료의 한 조각에서 $\Delta=\Delta_i>0$이라 하고, 기울기를 $\alpha=s_i/\Delta$, $\beta=s_{i+1}/\Delta$로 정규화하자. Fritsch와 Carlson은 조각의 3차 Hermite 다항식이 단조일 조건을 $(\alpha,\beta)$ 평면의 영역으로 완전히 기술했다\cite{fritsch1980}. 그 영역 안의 정사각형 하나는 손으로 확인할 수 있다.

\begin{proposition}
\label{prop:monotone-square}
$0\le\alpha,\beta\le3$이면 조각의 3차 Hermite 보간은 단조 증가한다.
\end{proposition}

\begin{proof}
척도를 바꾸어 조각을 $[0,1]$, 값을 $0$에서 $1$로 두면 $\Delta=1$이고, Hermite 기저의 도함수로 쓰면
\[
p'(u)=\alpha(1-u)(1-3u)+6u(1-u)+\beta u(3u-2)
\]
이다. 고정된 $u$에서 이것은 $(\alpha,\beta)$의 일차식이므로, 정사각형 위의 최솟값은 꼭짓점에서 얻어진다. 네 꼭짓점에서 차례로 $p'=6u(1-u)$, $3(1-u)^2$, $3u^2$, $3(1-2u)^2$이고 모두 $[0,1]$에서 $0$ 이상이다.
\end{proof}

\subsection{PCHIP and Akima}

\term{PCHIP}{piecewise cubic Hermite interpolating polynomial}는 이 조건을 만족하도록 기울기를 국소적으로 정한다. 이웃한 두 기울기 $\Delta_{i-1}$, $\Delta_i$의 부호가 다르거나 하나가 $0$이면 그 노드는 극값이어야 하므로 $s_i=0$으로 둔다. 같은 부호이면 Fritsch와 Butland의 가중 조화 평균
\[
\frac1{s_i}=\frac1{w_1+w_2}\Bigl(\frac{w_1}{\Delta_{i-1}}+\frac{w_2}{\Delta_i}\Bigr),\qquad w_1=2h_i+h_{i-1},\quad w_2=h_i+2h_{i-1}
\]
을 쓴다\cite{fritsch1984}. 분모에서 $w_2/\Delta_i$를 버리면 $s_i\le(1+w_2/w_1)\Delta_{i-1}$이고, $w_2/w_1\le2$이므로 $s_i\le3\Delta_{i-1}$이다. 같은 방법으로 $s_i\le3\Delta_i$이다. 그래서 정규화한 기울기가 \cref{prop:monotone-square}의 정사각형 안에 들어가고, PCHIP는 단조로운 자료를 단조롭게 보간한다. 앞의 계단 자료에서 PCHIP의 곡선은 정확히 $[0,1]$ 안에 머문다. 기울기 하나가 이웃한 두 조각만으로 정해지므로 계산이 국소적이고 병렬로 할 수 있다.

Akima는 가운데 두 차분 기울기 $m_{i-1}$, $m_i$를 무게 $|m_{i+1}-m_i|$, $|m_{i-1}-m_{i-2}|$로 평균해 기울기를 정하는 다른 국소 방법을 제안했다\cite{akima1970}. 무게는 차분 기울기의 변화량이어서 곡선이 꺾이는 곳 근처에서 꺾이지 않은 쪽의 기울기를 따른다. 단조성은 보장하지 않지만 PCHIP보다 덜 평평하고 출렁임은 3차 스플라인보다 적다. 평평한 구간이나 같은 값이 이어지는 자료에서 생기는 문제를 고친 판을 수정 Akima(makima)라 부른다.

\section{B-Splines}
\label{sec:bspline}

\subsection{Definition}

3차 스플라인을 기울기라는 노드의 양으로 정했지만 스플라인 전체가 이루는 함수 공간에는 국소적인 기저가 있다. 감소하지 않는 \term{매듭}{knot} $t_0\le t_1\le\dots$에 대해 차수 $0$의 B-스플라인을 $B_{j,0}(x)=1$ $(t_j\le x<t_{j+1})$, 나머지에서 $0$으로 두고
\begin{equation}
\label{eq:cox-de-boor}
B_{j,k}(x)=\frac{x-t_j}{t_{j+k}-t_j}B_{j,k-1}(x)+\frac{t_{j+k+1}-x}{t_{j+k+1}-t_{j+1}}B_{j+1,k-1}(x)
\end{equation}
로 차수를 올린 것이 차수 $k$의 \term{B-스플라인}{B-spline}이다($0/0$은 $0$으로 둔다). 이 점화식은 Cox와 de Boor가 각각 1972년에 내놓았다\cite{cox1972,deboor1972}. 스플라인이라는 이름과 B-스플라인의 기본 성질은 Schoenberg의 1946년 연구와 Curry와의 공동 연구에서 왔다\cite{schoenberg1946,curry1966}.

\begin{proposition}
\label{prop:bspline-properties}
$B_{j,k}$는 다음을 만족한다.
\begin{enumerate}
\item 받침이 $[t_j,t_{j+k+1})$이다. 곧 그 밖에서 $0$이다.
\item $B_{j,k}\ge0$이다.
\item 매듭 구간 $[t_l,t_{l+1})$에서 $0$이 아닌 것은 $B_{l-k,k},\dots,B_{l,k}$의 $k+1$개이고, 그 합은 $1$이다.
\end{enumerate}
\end{proposition}

\begin{proof}
(i)과 (ii)는 $k$에 대한 귀납법이다. 받침 $[t_j,t_{j+k})$와 $[t_{j+1},t_{j+k+1})$의 합집합이 $[t_j,t_{j+k+1})$이고, 그 위에서 \eqref{eq:cox-de-boor}의 두 계수 $(x-t_j)/(t_{j+k}-t_j)$와 $(t_{j+k+1}-x)/(t_{j+k+1}-t_{j+1})$은 해당 받침 위에서 $0$과 $1$ 사이에 있다. (iii)의 앞부분은 (i)에서 나온다. 합이 $1$이라는 것도 귀납법이다. $\sum_jB_{j,k}$에 \eqref{eq:cox-de-boor}를 넣고 같은 $B_{j,k-1}$의 계수를 모으면, $B_{j,k-1}$에는 $(x-t_j)/(t_{j+k}-t_j)$와 앞의 둘째 항에서 온 $(t_{j+k}-x)/(t_{j+k}-t_j)$가 붙어 합이 $1$이다. 그러므로 $\sum_jB_{j,k}=\sum_jB_{j,k-1}=\dots=\sum_jB_{j,0}=1$이다.
\end{proof}

매듭이 $m$겹인 자리에서 $k-m$번 연속 미분 가능한 차수 $k$의 스플라인을 모두 모은 공간을 생각하자. Curry와 Schoenberg는 이 공간이 B-스플라인들로 펼쳐지고 B-스플라인들이 그 기저를 이룸을 보였다\cite{curry1966}. 매듭을 겹치면 그 자리의 매끄러움을 줄일 수 있고 양 끝 매듭을 $k+1$겹으로 두면 끝점에서 보간하는 곡선이 된다. 스플라인을 $s(x)=\sum_jc_jB_{j,k}(x)$로 쓰는 것을 B-스플라인 표현이라 하고, $c_j$를 \term{조정점}{control point}이라 부르기도 한다. \cref{prop:bspline-properties}에 따라 $s(x)$는 이웃한 $k+1$개의 계수의 볼록 결합이므로 계수가 곡선의 모양을 국소적으로, 직관적으로 조정한다. 그래서 컴퓨터 그래픽스는 곡선과 곡면을 B-스플라인으로 그린다.

\subsection{Evaluation}

$x$에서 $s(x)$를 구하는 방법은 두 가지이다. 하나는 \eqref{eq:cox-de-boor}로 $0$이 아닌 B-스플라인 값 $k+1$개를 차수 $0$부터 올려 계수와 곱하는 것이다. 다른 하나는 계수 쪽에 같은 점화식을 쓰는 \term{de Boor 알고리즘}{de Boor's algorithm}이다. 구간 $[t_l,t_{l+1})$에서 $d_j^{(0)}=c_j$ $(j=l-k,\dots,l)$로 두고 $r=1,\dots,k$에 대해
\[
d_j^{(r)}=(1-\alpha_{j,r})\,d_{j-1}^{(r-1)}+\alpha_{j,r}\,d_j^{(r-1)},\qquad\alpha_{j,r}=\frac{x-t_j}{t_{j+k+1-r}-t_j}
\]
을 $j=l-k+r,\dots,l$에 대해 하면 $s(x)=d_l^{(k)}$이다. 어느 쪽이든 모든 단계가 $[0,1]$의 무게로 만든 볼록 결합이므로 반올림 오차가 커지지 않고 값이 넘치지도 않는다. 같은 볼록 결합 구조는 Boehm의 \term{매듭 삽입}{knot insertion}에도 쓰인다\cite{boehm1980}. 매듭을 하나 끼워 넣어 같은 곡선을 더 많은 계수로 나타내는 방법이다.

\subsection{Interpolation and Total Positivity}

점 $\tau_1<\dots<\tau_N$에서 값을 맞추는 스플라인 $\sum_jc_jB_j$의 계수는 \term{배치 행렬}{collocation matrix} $A_{ij}=B_j(\tau_i)$에 대한 계 $Ac=y$의 해이다. 각 행에 $0$이 아닌 원소가 $k+1$개 이하이므로 $A$는 띠 행렬이다.

\begin{theorem}[Schoenberg--Whitney]
\label{thm:schoenberg-whitney}
배치 행렬 $A$가 정칙일 필요충분조건은 모든 $i$에서 $B_i(\tau_i)\ne0$, 곧 각 점이 대응하는 B-스플라인의 받침 안에 있는 것이다.
\end{theorem}

이 정리는 Schoenberg와 Whitney가 증명했다\cite{schoenberg1953whitney}. 더 나아가 배치 행렬은 모든 소행렬식이 $0$ 이상인 \term{완전 비음}{totally nonnegative} 행렬이다\cite{karlin1968,deboor2001guide}. 완전 비음인 정칙 행렬은 피벗 없는 Gauss 소거의 $L$과 $U$가 모두 비음이라서 \cref{prop:thomas-positive} 뒤의 논증과 똑같이 $|L||U|=|A|$이고 성장이 없다. de Boor와 Pinkus는 이것으로 완전 비음 계를 피벗 없이 푸는 것이 성분별로 후진 안정함을 보였다\cite{deboor1977pinkus}. 부분 피벗은 오히려 이 구조를 깨뜨릴 수 있다.

\section{Splines in stardust}
\label{sec:spline-stardust}

\begin{instardust}[\code{docs/design/interpolate.md} §3.3]
stardust의 \code{CubicSpline}, \code{PchipInterpolator}, \code{Akima1DInterpolator}는 모두 같은 함수로 계수를 만든다. 노드의 기울기를 구한 뒤 \eqref{eq:hermite-coef}를 적용하는 함수이며 계수의 연산 순서는 scipy와 같다. \code{CubicSpline}의 삼중대각 계는 \cref{prop:thomas-positive}에 따라 피벗 없는 Thomas 소거로 푼다. 설계 문서는 not-a-knot과 주기 조건까지 넣어 피벗이 양수임을 증명하고, 계산한 기울기의 성분별 오차 상계를 정확도 계약으로 삼는다. scipy는 같은 계를 LAPACK의 부분 피벗 띠 풀이 \code{gbsv}로 푼다. 설계 문서의 측정에서는 scipy의 기울기가 같은 상계를 넘어 최대 $6.5$배에 이른 시험 곡선이 있었다. 기울기가 $0$인 평평한 걸음 근처에서 부분 피벗이 원래의 비음 구조를 깨뜨린 것이다.
\end{instardust}

\begin{instardust}[\code{docs/design/interpolate.md} §3.4]
B-스플라인의 평가는 \eqref{eq:cox-de-boor}를 차수 $0$부터 올리는데 두 몫 $(x-t_j)/(t_{j+k}-t_j)$를 먼저 만들어 곱한다. 두 몫이 모두 $[0,1]$이라 모든 단계가 넘치지 않는 볼록 결합이다. scipy는 B 값을 매듭 간격으로 먼저 나누므로 간격이 비정상 수일 만큼 좁으면 그 몫이 넘쳐 무한대가 된다. 보간의 배치 행렬은 도함수 조건이 없거나 끝에서 낮은 차수부터 차례로 주어지면 완전 비음 구조를 지키도록 행을 늘어놓고 피벗 없이 소거한다. 자연 경계 조건의 3차 스플라인에서는 피벗 하나가 음수가 된다. 설계 문서는 그 자리의 성장이 $2$ 미만임을 증명했고 정확한 유리수로 잰 값도 $1.99999\cdots$였다. 그 밖의 도함수 조건에서는 성장이 매듭 간격의 제곱에 반비례해 커질 수 있으므로 처음부터 부분 피벗으로 분해한다.
\end{instardust}

\section{Concluding Remarks}

스플라인은 차수를 낮게 묶어 두고 조각을 늘려 정확도를 얻는다. 그래서 노드를 고를 수 없는 자료에서도 안정하고, 계산이 국소적이며, 풀어야 할 계가 띠 행렬이다. 3차 스플라인은 곡률 에너지를 최소화하는 가장 매끄러운 보간이지만 그 매끄러움 때문에 자료 밖으로 출렁이고, PCHIP는 2차 도함수의 연속을 포기하고 모양을 지킨다. B-스플라인은 같은 공간에 국소적이고 비음인 기저를 주어 계수가 곡선을 볼록 결합으로 조정하게 하고 보간 행렬을 완전 비음으로 만든다. 다음 장에서는 변수가 여럿인 자료를 다룬다. 격자 위의 자료에는 이 장의 방법을 축마다 쓰는 텐서곱이 있고 흩어진 자료에는 전혀 다른 도구가 필요하다.

스플라인의 표준 교과서는 de Boor의 책\cite{deboor2001guide}이고, 이론적인 측면은 Schumaker의 책\cite{schumaker2007}이 자세하다.

\chapterreferences
